\section{Related Work}
\label{sec:related}

\paragraph{Informal precedents.}
Pairing probability with intuitionistic logic is not a new idea,
though the literature on it is scattered.  Weatherson develops
intuitionistic probability axiomatically and defends it with
Dutch-book arguments~\cite{weatherson2003}.  His axioms match our
valuation axioms exactly, with
negation derived rather than assumed, so the object itself and the
bound $v(a) + v(\neg a) \le 1$ are his.  What is new here
is the decomposition of that bound's slack into $\mathrm{dnGap}$ and
$\mathrm{deMorganGap}$ and its determinacy from $v$ alone
(Theorem~\ref{thm:slack-explicit}).  
Paris shows that the assignments avoiding a Dutch book are the mixtures
of a logic's two-valued valuations, and that when the valuations respect
$\wedge$ and $\vee$, as in intuitionistic and modal logics, they are
exactly the normalized, monotone, modular
assignments~\cite{paris2001}. On finite frames this is
Theorem~\ref{thm:mixture}. Williams obtains the same set from accuracy
domination, under the name generalized probabilism~\cite{williams2012}.
Paris also shows, from Shafer's results, that the mixtures of the
``known true'' valuations, which make a sentence true when it follows
from what is known, are exactly the Dempster--Shafer belief
functions~\cite{paris1994,paris2001,shafer1976}. These valuations do not
respect $\vee$.  We mechanize an analogous fact for our
intuitionistic setting: the $2$-monotone case follows directly from
modularity\leanref{two-monotone}, and the full tower is
Theorem~\ref{thm:infty}.  The
identification $\Bel = P(\Box\,\cdot)$ connecting belief
functions~\cite{dempster1967,shafer1976} to modal logic is
established~\cite{ruspini1986}.  Section~\ref{sec:bridges} shows this
identification is not a separate axiom for us but a special case of
our frame-level slack, recovered automatically once the frame is a
topology and $\neg$ is read as the interior of the complement
\leanref{tovaluationopens-slack-eq-frontier}.  

On the localic
side, valuations on
frames are standard objects of constructive measure theory, with
Vickers's work and the Coquand--Spitters development as canonical
references~\cite{vickers2008,coquand-spitters2009}.  These normally
build in Scott continuity and target additive, Riesz-style goals.
Section~\ref{sec:representation} shows Scott continuity to be a
substantive assumption (it decides point-representability on the
frames we examine), so our base structure omits it.  Cox, Jaynes,
Halpern's counterexample, and the Van Horn and Arnborg--Sj\"odin
repairs~\cite{cox1946,jaynes2003,halpern1999,vanhorn2003,
arnborg-sjodin2000} are classical throughout, and Cox's and Van
Horn's derivations both lean on double-negation elimination.
Section~\ref{sec:cox} proves
the sum-rule half without that move, at the price of positing
modularity as an axiom rather than deriving it
(Example~\ref{thm:no-functional}).  Van Horn observes
that dropping axiom R3 yields an interval-valued theory, exactly the
belief/plausibility pair our slack separates.  To our knowledge, no
prior work, informal or formal, states either characterization
theorem: complement rule $\iff$ excluded middle
(Theorem~\ref{thm:hinge}), or Dempster $=$ Bayes $\iff$ zero slack
at the evidence (Theorem~\ref{thm:cond-hinge}).

\paragraph{Prior formalizations.}
Two bodies of mechanized work are adjacent, and neither overlaps.
ALEA~\cite{audebaud-paulin2009} formalizes a discrete distribution
monad in Coq for randomized programs: additive probability, no
event logic in our sense.  Closest is the
HoTT/Coq line of Faissole and Spitters on \emph{valuations} over
lower reals~\cite{faissole-spitters2017}, done constructively with
Scott continuity and aimed at integration (Riesz) and a Giry-style
valuation \emph{monad}.  Their valuations differ from ours in
carrying Scott continuity, and their monad is the Scott-continuous
instance of the question examined in
Section~\ref{sec:monad}.  Our development drops the continuity
their base structure builds in, and the two mechanizations do not
otherwise overlap.  On the domain-theoretic side, commutativity of
the full valuations monad on the category of directed-complete
partial orders is open.  A hierarchy of commutative submonads,
generated from simple valuations (finite mixtures of point
valuations), has been constructed
instead~\cite{jia2021commutative}.  
Our countably-presented fragment, the valuations given as countable
mixtures (Proposition~\ref{prop:mixcountable}), is the frame-side
analogue.

\paragraph{Non-Boolean probability elsewhere.}
Quantum probability on orthomodular lattices, with Gleason's
theorem~\cite{birkhoff-vonneumann1936,gleason1957}, is the
long-standing precedent for re-basing probability on a non-Boolean
event logic.  Orthomodular lattices give up distributivity itself,
while our frames stay distributive by definition
(Definition~\ref{def:frame}), locating all the non-classical
behavior in the failure of excluded middle alone.
Section~\ref{sec:quantum} applies the construction from certain states
to a quantum event logic.  



Linear logic
marks a different kind of departure.  Its event algebras are
quantales, complete lattices equipped with an associative tensor
$\otimes$ that need not be idempotent~\cite{girard1987}.  The
lattice reduct of a quantale is already an instance of
Section~\ref{sec:lattice}'s generalization, but $\otimes$ itself is
not $\sqcap$, and a valuation grading it cannot be phrased in
Definition~\ref{def:lvaluation}'s vocabulary without first deciding
what modularity even means for a non-idempotent operation.
Section~\ref{sec:recipe} replaces modularity by the affine laws valid on
the certain states.  The probabilistic semantics that does exist for
linear logic, Girard's probabilistic coherence spaces as developed by
Danos and Ehrhard~\cite{danos-ehrhard2011}, grades the multiplicative
structure directly rather than valuing a lattice.
Section~\ref{sec:recipe} shows that it contains the calculus of
certain states at every multiplicative--additive formula and equals it
at implications between data types, and Section~\ref{sec:second} that
the inclusion is strict at second order.  



Free probability,
fuzzy measures, and imprecise
probability~\cite{walley1991} vary other parameters (independence,
additivity, single-valuedness).  Where Walley's coherent lower
previsions posit the belief/plausibility gap directly as a
coherence primitive, our slack (Section~\ref{sec:valuation}) is
derived, forced by modularity and negation rather than assumed.
Fritz's synthetic Markov-category
framework~\cite{fritz2020markov} varies the ambient structure
categorically rather than order-theoretically, axiomatizing
conditional independence and sufficient statistics abstractly
across discrete, measure-theoretic, and Gaussian probability alike.
It remains unmechanized and is orthogonal to our variation in the
underlying logic.  Wolpert and Kinney model
mathematics itself as a stochastic process, replacing binary
theoremhood with a probability distribution~\cite{wolpert-kinney2024}.
Theoremhood is a semidecidable event, so the computability guard of
Section~\ref{sec:bridges} bears directly on that program:
sharp credences about such events are uncomputable.  Our
variation is specifically the
\emph{logic}, holding the Cox-style methodology fixed.  Within DS
theory, the two conditionings and their inequivalence are
classical~\cite{shafer1976,fagin-halpern1991}.  What is new in
Theorem~\ref{thm:cond-hinge} is the exact logical characterization
of their agreement, and its mechanization.
